\chapter{THE BALANCING ROBOT MODEL}
\label{ch:model}

\section{MODEL OF A DIRECT CURRENT MOTOR}
\label{sec:motor}

The armature circuit of a brushed DC motor with terminal voltage $V_a$,
resistance $R$, inductance $L$, and back-EMF constant $k_e$ obeys Kirchhoff's
voltage law \cite{desilva2004mechatronics},
\begin{equation}\label{eq:circuit}
L\frac{\dd i}{\dd t}+Ri+k_e\omega=V_a ,
\end{equation}
where $i$ is the armature current and $\omega=\dot\vartheta$ the shaft rate.
Rearranging,
\begin{equation}\label{eq:circuit2}
\frac{\dd i}{\dd t}=\frac{V_a}{L}-\frac{R}{L}i-\frac{k_e}{L}\omega .
\end{equation}
The electromagnetic torque developed is proportional to the current,
\begin{equation}\label{eq:torque}
\tau_m=k_mi ,
\end{equation}
and the viscous friction torque is proportional to the shaft rate,
\begin{equation}\label{eq:friction}
\tau_f=k_f\omega .
\end{equation}
Applying Newton's second law and summing moments about the motor output shaft
against an applied load torque $\tau_a$ and load inertia $I_m$,
\begin{equation}\label{eq:shaft}
I_m\dot\omega=\tau_m-\tau_f-\tau_a .
\end{equation}
Substituting \eqref{eq:torque} and \eqref{eq:friction} into \eqref{eq:shaft} and
rearranging gives the second equation of motion,
\begin{equation}\label{eq:shaft2}
\frac{\dd\omega}{\dd t}=\frac{k_m}{I_m}i-\frac{k_f}{I_m}\omega-\frac{\tau_a}{I_m} .
\end{equation}

Equations \eqref{eq:circuit2} and \eqref{eq:shaft2} describe the fundamental
behavior of a DC motor as a pair of first-order differential equations relating
circuit current, motor speed, and applied torque. For a balancing robot a
simplified model provides acceptable performance. Assuming the motor inductance
$L$ and the friction coefficient $k_f$ are negligible,
\begin{equation}\label{eq:isimple}
i=\frac{V_a}{R}-\frac{k_e}{R}\omega ,
\end{equation}
and
\begin{equation}\label{eq:wsimple}
\frac{\dd\omega}{\dd t}=\frac{k_m}{I_m}i-\frac{\tau_a}{I_m} .
\end{equation}
Substituting \eqref{eq:isimple} into \eqref{eq:wsimple} gives a model that does
not include the circuit current,
\begin{equation}\label{eq:motor}
\frac{\dd\omega}{\dd t}
=-\frac{k_mk_e}{I_mR}\omega+\frac{k_m}{I_mR}V_a-\frac{\tau_a}{I_m} .
\end{equation}
This elimination is desirable: the current is not directly commandable on the
physical system and would add a state with no observable benefit. Neglecting the
inductance amounts to assuming the current reaches steady state immediately
relative to the motor speed, which is justified because the electrical time
constant $L/R$ is orders of magnitude shorter than the mechanical one. In
state-space form,
\begin{equation}\label{eq:motorss}
\frac{\dd}{\dd t}
\begin{bmatrix}\vartheta\\\omega\end{bmatrix}
=\begin{bmatrix}
0 & 1\\[4pt]
0 & -\dfrac{k_mk_e}{I_mR}
\end{bmatrix}
\begin{bmatrix}\vartheta\\\omega\end{bmatrix}
+\begin{bmatrix}
0 & 0\\[4pt]
\dfrac{k_m}{I_mR} & -\dfrac{1}{I_m}
\end{bmatrix}
\begin{bmatrix}V_a\\\tau_a\end{bmatrix}.
\end{equation}

\section{TWO WHEELED INVERTED PENDULUM DYNAMIC MODEL}
\label{sec:twipmodel}

The TWIP shares many similarities with the pendulum-on-a-cart problem. The
motion is comparable but the dynamics are more involved. The analysis begins
with the two wheels, the pendulum is analyzed separately, and the two sets of
equations are combined by eliminating the interface reactions, yielding two
dynamic equations describing the motion of the system.

\subsection{Wheel Dynamics}

The free-body diagram of the left wheel is shown in Figure~\ref{fig:fbdwheel}.
The reaction forces between the wheel and the pendulum are $P_L$ and $H_L$, the
torque applied to the wheel by the DC motor is $C_L$, and the friction force
between the ground and the wheel is $H_{fL}$. The equations for the left and
right wheels are completely analogous, so only the left wheel is analyzed. The
wheels are assumed to remain in ground contact at all times with no slip.

\begin{figure}[htbp]
\centering
\includegraphics[width=0.45\textwidth]{fbd_wheel}
\caption{Left wheel free body diagram.}
\label{fig:fbdwheel}
\end{figure}

Summing forces in the $x$ direction,
\begin{equation}\label{eq:wheelx}
M_w\ddot x=H_{fL}-H_L .
\end{equation}
Summing moments about the center of the wheel,
\begin{equation}\label{eq:wheelm}
I_w\ddot\theta_w=C_L-H_{fL}r ,
\end{equation}
where $I_w$ is the wheel moment of inertia and $\ddot\theta_w$ the wheel angular
acceleration. Starting from the motor shaft equation \eqref{eq:shaft} and
substituting \eqref{eq:isimple}, the motor torque delivered to either wheel is
\begin{equation}\label{eq:CL}
C_L=C_R=\tau_m=-\frac{k_mk_e}{R}\dot\theta_w+\frac{k_m}{R}V_a .
\end{equation}
\rev{In \eqref{eq:CL}, $\dot\theta_w$ must be read as the wheel rate relative to
the motor stator, which is fixed to the pendulum body. The development below
follows revision~7 in substituting the absolute rate $\dot x/r$, which drops a
tilt-rate term; the correction used for the results is given in
Section~\ref{sec:simmodel}.}
Using \eqref{eq:CL} in \eqref{eq:wheelm} and rearranging gives the friction
reaction force,
\begin{equation}\label{eq:HfL}
H_{fL}=-\frac{k_mk_e}{Rr}\dot\theta_w+\frac{k_m}{Rr}V_a
-\frac{I_w}{r}\ddot\theta_w .
\end{equation}
Substituting \eqref{eq:HfL} into \eqref{eq:wheelx} gives the equation of motion
for the left wheel,
\begin{equation}\label{eq:wheelL}
M_w\ddot x=-\frac{k_mk_e}{Rr}\dot\theta_w+\frac{k_m}{Rr}V
-\frac{I_w}{r}\ddot\theta_w-H_L ,
\end{equation}
and an identical development for the right wheel gives
\begin{equation}\label{eq:wheelR}
M_w\ddot x=-\frac{k_mk_e}{Rr}\dot\theta_w+\frac{k_m}{Rr}V
-\frac{I_w}{r}\ddot\theta_w-H_R .
\end{equation}

The no-slip assumption provides the transformation between rotational and
linear motion,
\begin{equation}\label{eq:noslip}
\ddot\theta_w=\frac{\ddot x}{r},\qquad \dot\theta_w=\frac{\dot x}{r} .
\end{equation}
Applying \eqref{eq:noslip} to \eqref{eq:wheelL} and \eqref{eq:wheelR},
\begin{align}
M_w\ddot x&=-\frac{k_mk_e}{Rr^2}\dot x+\frac{k_m}{Rr}V
-\frac{I_w}{r^2}\ddot x-H_L ,\label{eq:wheelL2}\\[4pt]
M_w\ddot x&=-\frac{k_mk_e}{Rr^2}\dot x+\frac{k_m}{Rr}V
-\frac{I_w}{r^2}\ddot x-H_R .\label{eq:wheelR2}
\end{align}
Adding \eqref{eq:wheelL2} and \eqref{eq:wheelR2} gives the equation of motion
for the combined wheel system,
\begin{equation}\label{eq:wheelcomb}
2\left(M_w+\frac{I_w}{r^2}\right)\ddot x
=-\frac{2k_mk_e}{Rr^2}\dot x+\frac{2k_m}{Rr}V-\bigl(H_L+H_R\bigr).
\end{equation}

\subsection{Pendulum Dynamics}

Figure~\ref{fig:fbdpend} shows the free-body diagram of the inverted pendulum.
Two reference frames are used: an inertial frame attached to the base of the
pendulum, and a translating frame at the center of gravity. Here $l$ is the
distance from the base to the center of gravity; $\alpha_p$, $\omega_p$, and
$\theta$ are the angular acceleration, angular velocity, and angular position of
the pendulum; $x,\dot x,\ddot x$ are the position, velocity, and acceleration of
the pendulum base; and $x_p,\dot x_p,\ddot x_p$ those of the center of gravity.

\begin{figure}[htbp]
\centering
\includegraphics[width=0.45\textwidth]{fbd_pendulum}
\caption{Inverted pendulum free body diagram.}
\label{fig:fbdpend}
\end{figure}

Summing forces in the $x_p$ direction,
\begin{equation}\label{eq:pendx}
\sum F_{xp}=H_L+H_R=M_p\ddot x_p ,
\end{equation}
and in the $y_p$ direction,
\begin{equation}\label{eq:pendy}
\sum F_{yp}=P_L+P_R-M_pg=M_p\ddot y_p .
\end{equation}

The acceleration of a point in a translating frame follows from relative motion
analysis \cite{hibbeler2010dynamics}. The acceleration of the pendulum center of
gravity, denoted $\mathbf{a}_B$, is
\begin{equation}\label{eq:relmotion}
\mathbf{a}_B=\mathbf{a}_A
+\bm{\alpha}_p\times\mathbf{r}_{B/A}
+\bm{\omega}_p\times\bigl(\bm{\omega}_p\times\mathbf{r}_{B/A}\bigr),
\end{equation}
with
\begin{align}
\mathbf{a}_B&=\ddot x_p\hat\imath+\ddot y_p\hat\jmath ,\label{eq:aB}\\
\mathbf{r}_{B/A}&=l\sin\theta\,\hat\imath-l\cos\theta\,\hat\jmath ,\label{eq:rBA}\\
\bm{\omega}_p&=\dot\theta\,\hat k ,\label{eq:omegap}\\
\bm{\alpha}_p&=\ddot\theta\,\hat k ,\label{eq:alphap}\\
\mathbf{a}_A&=\ddot x\,\hat\imath .\label{eq:aA}
\end{align}
Using \eqref{eq:rBA}--\eqref{eq:aA} in \eqref{eq:relmotion},
\begin{equation}\label{eq:aBfull}
\mathbf{a}_B=\Bigl(\ddot x+\ddot\theta l\cos\theta-\dot\theta^2l\sin\theta\Bigr)\hat\imath
+\Bigl(\ddot\theta l\sin\theta+\dot\theta^2l\cos\theta\Bigr)\hat\jmath .
\end{equation}
Substituting \eqref{eq:aBfull} into \eqref{eq:pendx} and \eqref{eq:pendy},
\begin{align}
H_L+H_R&=M_p\Bigl(\ddot x+\ddot\theta l\cos\theta-\dot\theta^2l\sin\theta\Bigr),
\label{eq:Hsum}\\[4pt]
P_L+P_R&=M_pg+M_p\Bigl(\ddot\theta l\sin\theta+\dot\theta^2l\cos\theta\Bigr).
\label{eq:Psum}
\end{align}

Summing moments about the pendulum center of gravity, with the clockwise
direction taken as positive,
\begin{equation}\label{eq:pendm}
\sum M=-\bigl(C_L+C_R\bigr)-\bigl(H_L+H_R\bigr)l\cos\theta
-\bigl(P_L+P_R\bigr)l\sin\theta=I_p\ddot\theta .
\end{equation}

\subsection{Combined Equations of Motion}

Combining \eqref{eq:CL}, \eqref{eq:Hsum}, \eqref{eq:Psum}, and \eqref{eq:pendm},
\begin{equation}\label{eq:combined}
\begin{aligned}
-\left(-\frac{2k_mk_e}{R}\dot\theta_w+\frac{2k_m}{R}V_a\right)
&-\Bigl(M_p\bigl(\ddot x+\ddot\theta l\cos\theta-\dot\theta^2l\sin\theta\bigr)\Bigr)l\cos\theta\\
&-\Bigl(M_pg+M_p\bigl(\ddot\theta l\sin\theta+\dot\theta^2l\cos\theta\bigr)\Bigr)l\sin\theta
=I_p\ddot\theta .
\end{aligned}
\end{equation}
Cancelling terms, rearranging, and applying the identity
$\sin^2\theta+\cos^2\theta=1$ gives the first nonlinear equation of motion,
\begin{equation}\label{eq:nl1}
\bigl(I_p+M_pl^2\bigr)\ddot\theta
-\frac{2k_mk_e}{Rr}\dot x
+\frac{2k_m}{R}V
+M_pgl\sin\theta
=-M_pl\ddot x\cos\theta .
\end{equation}
Using \eqref{eq:Hsum} to remove the body forces from \eqref{eq:wheelcomb} gives
the second nonlinear equation of motion,
\begin{equation}\label{eq:nl2}
\frac{2k_m}{Rr}V=\left(2M_w+\frac{2I_w}{r^2}+M_p\right)\ddot x
+\frac{2k_mk_e}{Rr^2}\dot x
+M_pl\ddot\theta\cos\theta
-M_pl\dot\theta^2\sin\theta .
\end{equation}

Equations \eqref{eq:nl1} and \eqref{eq:nl2} describe the full nonlinear motion
of the two-wheeled inverted pendulum.

\subsection{Linearization}

In the convention of Figure~\ref{fig:fbdpend}, $\theta=\pi$ is the upright
vertical configuration. A set of linearized equations is obtained by writing
$\theta=\pi+\phi$, where $\phi$ is a small angle from the vertical. The
small-angle assumptions are
\begin{equation}\label{eq:smallangle}
\cos\theta=-1,\qquad \sin\theta=-\phi,
\end{equation}
together with the assumption that the squared angular rate is small,
\begin{equation}\label{eq:smallrate}
\dot\theta^2=0 .
\end{equation}
Applying \eqref{eq:smallangle} and \eqref{eq:smallrate} to \eqref{eq:nl1} and
\eqref{eq:nl2} yields the linearized equations of motion,
\begin{align}
\bigl(I_p+M_pl^2\bigr)\ddot\phi-\frac{2k_mk_e}{Rr}\dot x
+\frac{2k_m}{R}V-M_pgl\phi&=M_pl\ddot x ,\label{eq:lin1}\\[4pt]
\frac{2k_m}{Rr}V=\left(2M_w+\frac{2I_w}{r^2}+M_p\right)\ddot x
+\frac{2k_mk_e}{Rr^2}\dot x-M_pl\ddot\phi&. \label{eq:lin2}
\end{align}

\rev{\begin{remark}[Sign of the gravity term]\label{rem:gravsign}
Rearranging \eqref{eq:lin1} gives
$\bigl(I_p+M_pl^2\bigr)\ddot\phi=M_pgl\phi+\cdots$, so the gravitational term
enters with positive sign and is destabilizing. This is the expected signature
of an upright linearization and confirms the convention: $\phi$ is the
displacement from the \emph{unstable} equilibrium, not from the hanging one. It
is also consistent with \eqref{eq:linss}, whose $(2,3)$ and $(4,3)$ entries are
both positive. The description of $\theta=\pi$ in revision~7 as the
``statically stable vertical point'' is a mislabeling; the corrected reading is
given in Section~\ref{sec:modelsim}.
\end{remark}}

Rearranging \eqref{eq:lin1} and \eqref{eq:lin2} into state-space form
$\dot{\mathbf{x}}=\mathcal{A}\mathbf{x}+\mathcal{B}V$ with
$\mathbf{x}=\begin{bmatrix}x&\dot x&\phi&\dot\phi\end{bmatrix}^T$ gives
\begin{equation}\label{eq:linss}
\mathcal{A}=
\begin{bmatrix}
0 & 1 & 0 & 0\\[4pt]
0 & \dfrac{2k_mk_e\bigl(M_plr-I_p-M_pl^2\bigr)}{Rr^2\alpha}
  & \dfrac{M_p^2gl^2}{\alpha} & 0\\[8pt]
0 & 0 & 0 & 1\\[4pt]
0 & \dfrac{2k_mk_e\bigl(r\beta-M_pl\bigr)}{Rr^2\alpha}
  & \dfrac{M_pgl\beta}{\alpha} & 0
\end{bmatrix},
\qquad
\mathcal{B}=
\begin{bmatrix}
0\\[4pt]
\dfrac{2k_m\bigl(I_p+M_pl^2-M_plr\bigr)}{Rr\alpha}\\[8pt]
0\\[4pt]
\dfrac{2k_m\bigl(M_pl-r\beta\bigr)}{Rr\alpha}
\end{bmatrix},
\end{equation}
where
\begin{equation}\label{eq:beta}
\beta:=2M_w+\frac{2I_w}{r^2}+M_p
\end{equation}
and
\begin{equation}\label{eq:alphadef}
\alpha:=I_p\beta+2M_pl^2\left(M_w+\frac{I_w}{r^2}\right).
\end{equation}

For the development in Chapter~\ref{ch:control} it is convenient to name the
entries of \eqref{eq:linss} directly. Writing $x_1=x$, $x_2=\dot x$, $x_3=\phi$,
$x_4=\dot\phi$, $u=V$, and appending the unmodeled and parametric uncertainty as
$d_1$ and $d_2$,
\begin{equation}\label{eq:plantform}
\begin{aligned}
\dot x_1&=x_2,\\
\dot x_2&=A_1x_2+A_2x_3+B_1u+d_1(\mathbf{x}),\\
\dot x_3&=x_4,\\
\dot x_4&=A_3x_2+A_4x_3+B_2u+d_2(\mathbf{x}),
\end{aligned}
\end{equation}
with $A_1,\dots,A_4,B_1,B_2$ read off from \eqref{eq:linss}. Note that
$A_4=0$ in the nominal model, since the tilt rate does not appear in the tilt
acceleration once friction is neglected; it is retained symbolically because the
uncertainty $d_2$ generally reintroduces such a dependence.

\rev{\begin{remark}[Structural consequences]
Two features of \eqref{eq:plantform} drive everything that follows. First,
$B_1$ and $B_2$ are both nonzero and the input is scalar, so the system is
underactuated and no choice of $u$ cancels both $d_1$ and $d_2$: the uncertainty
is unmatched. Second, $A_2\neq0$, so the tilt state enters the translational
acceleration directly. This is what makes $x_3$ usable as a virtual control for
the position subsystem, and it is also the term that couples the two error
subsystems in the backstepping design. It must be retained; see
Section~\ref{sec:cffb}.
\end{remark}}

\section{MODEL SIMULATION STUDY}
\label{sec:modelsim}

The model should be checked for accuracy and for the expected motion of an
inverted pendulum. This is done by examining the uncontrolled motion of the
nonlinear equations \eqref{eq:nl1}--\eqref{eq:nl2}; any consistent set of
parameter values may be used.

\rev{Two equilibria exist. In the convention of \eqref{eq:smallangle},
$\theta=\pi$ is the upright configuration and $\theta=0$ the hanging one. The
hanging point is statically stable; the upright point is an equilibrium but is
unstable, as Remark~\ref{rem:gravsign} establishes, so a trajectory started
exactly there with zero rate remains there while any perturbation grows.
Revision~7 described the two in the opposite sense. The numerical tests below
are unchanged; only their interpretation is corrected.}

Starting from $\theta=\SI{180}{\degree}$ with zero initial rate, the system
remains at that point, as shown in Figure~\ref{fig:stab180}. \rev{This confirms
that $\theta=\pi$ is an equilibrium of the implemented equations. It is not
evidence of stability, and a small perturbation about this point diverges, which
is the behavior required of an inverted pendulum and the reason active control
is needed at all.}

Starting from $\theta=\SI{10}{\degree}$, the system is expected to oscillate
about the hanging equilibrium $\theta=0$. With no friction modeled, the motor
back-EMF term supplies a damping-like effect and the oscillations decay.
Examining the free-body diagram in Figure~\ref{fig:fbdpend}, a small positive
angle should cause the body to move initially in the positive direction as the
tilt returns toward zero. Figure~\ref{fig:stab10} shows this behavior.
\rev{With the corrected back-EMF constant of Section~\ref{sec:simmodel}, the
damping is strong and the oscillation decays within about five periods. The
translational position does not return to zero; after the initial positive
excursion the base settles about \SI{8}{mm} behind its starting point, the
net effect of the asymmetric damping over the swings. The figures show the
corrected model used for the results of Chapter~\ref{ch:results}.} These checks
establish that the nonlinear equations of motion behave as expected and can be
relied upon.

\begin{figure}[htbp]
\centering
\includegraphics[width=0.8\textwidth]{stability_180}
\caption{Uncontrolled response from $\theta=\SI{180}{\degree}$.}
\label{fig:stab180}
\end{figure}

\begin{figure}[htbp]
\centering
\includegraphics[width=0.8\textwidth]{stability_10}
\caption{Uncontrolled response from $\theta=\SI{10}{\degree}$.}
\label{fig:stab10}
\end{figure}
